% Generated draft: requires mathematical and source review.
\documentclass[11pt]{article}
\usepackage[margin=25mm]{geometry}
\usepackage{fontspec}
\setmainfont{Latin Modern Roman}
\usepackage{amsmath,amssymb,mathtools,unicode-math}
\setmathfont{Latin Modern Math}
\usepackage{enumitem,needspace,xurl,hyperref}
\setlength{\parindent}{0pt}
\setlength{\parskip}{5pt}
\setlength{\emergencystretch}{3em}
\newcommand{\N}{\mathbb{N}}
\newcommand{\Z}{\mathbb{Z}}
\newcommand{\Q}{\mathbb{Q}}
\newcommand{\R}{\mathbb{R}}
\newcommand{\C}{\mathbb{C}}
\newcommand{\F}{\mathbb{F}}
\newcommand{\A}{\mathbb{A}}
\newcommand{\PP}{\mathbb{P}}
\newcommand{\E}{\mathbb{E}}
\newcommand{\eps}{\varepsilon}
\newcommand{\dd}{\,\mathrm{d}}
\begin{document}
\section*{UnsolvedMath ID: 30003191}
\textbf{Source problem code:} catalogue entry
\subsection{Definitions and mathematical statement}
For a prime number \(p\), write \(\Q_p\) for the completion of \(\Q\) with respect to the absolute value normalized by \(|p|_p=p^{-1}\). The source does not specify a particular prime or a field over \(\Q_p\).

To explain the classical terminology, a pure rational variation of Hodge structure of weight \(w\in\Z\) on a complex manifold \(S\) consists of a finite-rank rational local system \(\mathbb V\) and a decreasing holomorphic filtration \(F^\bullet\) by subbundles of \(\mathcal V=\mathbb V\otimes_{\Q}\mathcal O_S\). At each \(s\in S\), this filtration determines a decomposition
\[
\mathbb V_s\otimes_{\Q}\C=\bigoplus_{a+b=w}H_s^{a,b},\qquad
\overline{H_s^{a,b}}=H_s^{b,a},\qquad
F_s^r=\bigoplus_{a\ge r}H_s^{a,w-a}.
\]
The flat connection \(\nabla\) induced by \(\mathbb V\) satisfies Griffiths transversality:
\[
\nabla(F^r)\subseteq F^{r-1}\otimes_{\mathcal O_S}\Omega_S^1.
\]
After a local flat identification of the fibers, its period map sends \(s\) to the filtration \(F_s^\bullet\) in the corresponding space of Hodge filtrations. These definitions provide one classical reference setting; the source does not restrict its question to pure rational variations or specify polarization or additional structures.

\textbf{Source question.} What is the \(p\)-adic analogue of the classical theory of period maps and variations of Hodge structure?

This is a request for a theory, rather than a specified theorem or conjecture. The record does not prescribe a category of base spaces or families, coefficient fields, a notion of \(p\)-adic variation, a period target, or the classical properties that an analogue must reproduce. Consequently, it does not determine a unique mathematical assertion with explicit hypotheses and conclusion. No particular construction or completeness criterion is imposed here.
\subsection{Short English statement}
The problem asks for a \(p\)-adic counterpart of classical varying Hodge structures and their period maps. Its intended scope and success criteria are unspecified, so the supplied question remains a research program requiring clarification and human mathematical review.
\subsection{Sources}
\begin{itemize}
\item Supplied UnsolvedMath local registry record, permanent ID \(30003191\), title ``p-Adic Period Maps and Hodge-Theoretic Variations,'' source problem code \texttt{catalogue entry}, in the Oberwolfach Reports Open Problems collection. No individual proposer is supplied; the proposed year is recorded as 2016.
\item Supplied report citation: \emph{Arithmetic Geometry (2017)}, Oberwolfach Reports, DOI \texttt{10.4171/owr/2016/38}, pp.\ 2171--2224. \url{https://doi.org/10.4171/owr/2016/38}. The background describes these reports as documenting mathematical workshops and their open-problem discussions. The citation's year 2017 and the record's proposed year 2016 are retained without reconciliation.
\item Supplied literature reference: ``Scholze, p-adic period maps.'' \url{https://doi.org/10.4007/annals.2015.182.3.3}. This is the reference label supplied by the record; fuller bibliographic details are not supplied.
\item Supplied literature reference: ``Oberwolfach Report 38.'' \url{https://doi.org/10.4171/owr/2016/38}.
\item Literature assessment dated exactly \texttt{2026-08-21}: the supplied status is \texttt{partially\_solved}. The supplied assessment attributes broad \(p\)-adic Hodge--Tate period-map constructions to Scholze and subsequent work, while asserting that there is no single complete \(p\)-adic replacement for every classical variation-of-Hodge-structure feature. This is supplied dated curation derived from the OWR record, not an independently verified literature assessment; its time-sensitive status and cited claims require independent checking.
\item The registry labels the statement \texttt{corrected\_verified} and reports that the official report was checked and that this is its organizing open question for \(p\)-adic period maps. These are supplied verification claims, not verification performed for this draft.
\end{itemize}
\end{document}
